\chapter{振动控制与工程应用}
\intro{振动控制是振动理论从分析走向设计的关键环节。本章涵盖隔振、动力吸振、阻尼技术、主动控制及典型工程案例（高层建筑 TMD、转子动力学、汽车悬挂、桥梁风振）。}

\section{振动控制策略分类}
振动控制的目标是抑制有害振动、抑制振动产生的噪声、避免共振使设备处于安全工作范围。主要手段分为以下几类：

\subsection{消振（振源抑制）}
根本性地消除或削弱振动源头：平衡旋转部件以减小不平衡激振力、优化叶片设计降低气动激振、减小冲压车间锻锤的冲击力、采用减振刀杆降低切削颤振等。

\subsection{隔振（Vibration Isolation）}
将振源与对象通过弹性元件隔离，阻止振动能量传播。按隔离方向分为：
\begin{itemize}
    \item \textbf{力隔振（Force Isolation）}：减小振动机器传递至基础的动态力
    \item \textbf{基础隔振（Motion Isolation）}：减小基础运动传递给精密设备的振动
\end{itemize}
详见第 \S 2 节。

\subsection{动力吸振（Dynamic Vibration Absorber, DVA）}
在主系统上附加一个辅助质量-弹簧-阻尼子系统，利用反共振原理降低主系统在特定频率处的振幅。详见第 \S 3 节。

\subsection{阻尼减振}
利用阻尼材料的耗能特性将振动机械能转化为热能，尤其适用于共振峰值的削减和冲击响应的快速衰减。详见第 \S 4 节。

\subsection{主动/半主动控制}
利用传感器检测振动状态，通过控制器驱动作动器施加控制力，实现可调节的实时振动抑制。详见第 \S 5 节。

\section{隔振理论}
\subsection{力传递率（Force Transmissibility）}
考虑质量为 $m$ 的机器置于刚度 $k$、阻尼 $c$ 的隔振器上，机器受简谐激振力 $F_0\sin\omega t$ 的作用。传递到基础的动态力为：
\[
F_T(t) = k x(t) + c\dot{x}(t)
\]

\begin{myformula}
力传递率 $T_f$ 定义为传递力幅与激振力幅之比：
\[
T_f = \frac{|F_T|_{\max}}{F_0} = \frac{\sqrt{1 + [2\zeta(\omega/\omega_n)]^{2}}}{\sqrt{[1 - (\omega/\omega_n)^{2}]^{2} + [2\zeta(\omega/\omega_n)]^{2}}}
\]
其中 $\omega_n = \sqrt{k/m}$, $\zeta = c/(2m\omega_n)$。
\end{myformula}

\subsection{隔振效率与设计准则}
\begin{definition}[隔振的有效频率范围]
力传递率 $T_f < 1$（即有隔振效果）的条件为：
\[
\frac{\omega}{\omega_n} > \sqrt{2} \approx 1.414
\]
在此区域内，频率比越大（即隔振器越软）、阻尼越小，隔振效果越好。但过小的阻尼会恶化共振穿越时的性能。
\end{definition}

\begin{enumerate}
    \item \textbf{设计准则}：选择隔振器使系统固有频率远低于工作频率。一般满足 $\omega/\omega_n > 3 \sim 4$，传递率可降至 $10\%$ 以下
    \item \textbf{阻尼的权衡}：增大阻尼可抑制共振区的振幅峰值，但在隔振区会略微减小隔振效果
    \item \textbf{静态变形约束}：软弹簧虽提供好的高频隔振，但会引入较大的静态下沉量 $\delta_{\text{st}} = mg/k$
\end{enumerate}

\begin{figure}[H]
\centering
\begin{tikzpicture}[scale=1.0]
    \def\zvals{{0.05,0.1,0.2,0.5,1.0}}
    \draw[-Stealth] (0,0) -- (7,0) node[right] {$\omega/\omega_n$};
    \draw[-Stealth] (0,0) -- (0,4.5) node[above] {$T_f$};
    \draw[gray,dashed] (0,1) -- (7,1) node[right] {$T_f = 1$};
    \draw[gray,dashed] (2.0,0) -- (2.0,4.2);
    \node[gray] at (2.0,4.4) {$\sqrt{2}$};

    \foreach \z in {0.05,0.1,0.2,0.5,1.0} {
        \pgfmathsetmacro\idx{int(\z*100)}
        \draw[blue!\idx!red,thick] plot[domain=0.05:6.8,samples=200]
            (\x,{sqrt(1 + (2*\z*\x)^2) / sqrt((1-\x*\x)^2 + (2*\z*\x)^2)});
    }

    \node[blue] at (5.5,0.3) {$\zeta \to 0$};
    \node[red] at (6.2,1.2) {$\zeta=1.0$};

    \draw[green!50!black,fill=green!10] (2.0,0.5) rectangle (6.8,0.01);
    \node[green!50!black] at (4.5,0.2) {隔振区 $T_f<1$};
\end{tikzpicture}
\caption{力传递率曲线。红色为不同阻尼比对应的曲线。所有曲线在 $\omega/\omega_n = \sqrt{2}$ 处相交于 $T_f = 1$，右侧为隔振效果区。}
\label{fig:ch12_transmissibility}
\end{figure}

\subsection{基础隔振（位移传递率）}
基础作简谐运动 $y(t) = Y_0\sin\omega t$ 时，隔振对象的位移传递率与力传率形式相同：
\[
T_d = \frac{|X|_{\max}}{Y_0} = T_f
\]
精密设备的隔振设计应使 $\omega_n \ll \omega$，即在极软的隔振元件上安装质量块。

\subsection{隔振材料选择}
常用的隔振材料与元件：

\begin{itemize}
    \item \textbf{橡胶隔振器}：天然橡胶和合成橡胶，阻尼较大（$\zeta \approx 0.05 \sim 0.15$），刚度可通过配方和形状调节。工作频率范围 $5 \sim 30\,\text{Hz}$，适中小载荷
    \item \textbf{金属螺旋弹簧隔振器}：阻尼极小（$\zeta \approx 0.005$），大变形，承载能力强。常与橡胶垫串联增加阻尼
    \item \textbf{空气弹簧}：刚度极低（频率可降至 $1 \sim 2\,\text{Hz}$），承载能力可调，用于高级隔振台、车辆空气悬挂
    \item \textbf{钢丝网垫隔振器}：金属丝编织压缩成型，具有干摩擦阻尼，耐高温和腐蚀，适用于航空发动机、舰载设备等严酷环境
    \item \textbf{组合隔振器}：两段隔振（双层隔振）可提供更好的高频衰减斜率（$-40\,\text{dB/dec}$ 替代单层  ${-20}\,\text{dB/dec}$）
\end{itemize}

\section{动力吸振器设计}
\subsection{基本动力吸振器（无阻尼 DVA）}
在质量 $m_1$ 的主系统（刚度 $k_1$，受简谐力 $F_0\sin\omega t$ 激励）上附加无阻尼质量-弹簧副系统 $(m_2, k_2)$：

\begin{myformula}
耦合系统的运动方程：
\[
\begin{cases}
m_1\ddot{x}_1 + (k_1 + k_2)x_1 - k_2 x_2 = F_0\sin\omega t \\
m_2\ddot{x}_2 + k_2(x_2 - x_1) = 0
\end{cases}
\]
假设稳态解 $x_j = X_j \sin\omega t$，得主系统振幅：
\[
X_1 = \frac{(k_2 - m_2\omega^{2})F_0}{(k_1 + k_2 - m_1\omega^{2})(k_2 - m_2\omega^{2}) - k_2^{2}}
\]
当 $\omega^{2} = k_2/m_2 = \omega_a^{2}$（吸振器固有频率等于激励频率）时，分子为零，主系统振幅 $X_1 = 0$——即完全消振（反共振条件）。
\end{myformula}

\subsection{最优调谐比和最优阻尼比（有阻尼 DVA）}
实际 DVA 引入阻尼 $c_2$，形成 Voigt 型动力吸振器。考虑主系统无阻尼时（$\zeta_1 = 0$）的设计问题。定义参数：
\[
\mu = \frac{m_2}{m_1}, \quad \alpha = \frac{\omega_a}{\omega_1}, \quad f = \frac{\omega}{\omega_1}, \quad \zeta_a = \frac{c_2}{2m_2\omega_2}
\]

\begin{myformula}
经典 Den Hartog 最优调谐公式（$\zeta_1 = 0$，$H_\infty$ 优化）：
\[
\alpha_{\text{opt}} = \frac{1}{1+\mu}
\]
\[
\zeta_{\text{opt}} = \sqrt{\frac{3\mu}{8(1+\mu)}}
\]
替代的 $H_2$ 优化公式（针对宽带激励的均方响应最小化）：
\[
\alpha_{\text{opt}} = \frac{1}{\sqrt{1+\mu}}, \quad \zeta_{\text{opt}} = \sqrt{\frac{\mu(4+3\mu)}{8(1+\mu)(2+\mu)}}
\]
\end{myformula}

图~\ref{fig:ch12_tmd} 展示了 TMD 调谐质量阻尼器的机械模型。

\begin{figure}[H]
\centering
\begin{tikzpicture}[scale=1.1]
    \draw[vibWall] (-3,-3) rectangle (-2.5,3);

    \draw[gray!30] (-2.5,0) -- (2.5,0);
    \draw[pattern=north east lines] (-1.0,-0.2) rectangle (1.0,0);

    \draw[vibSpring] (-2.5,1.0) -- (-0.5,1.0);
    \draw[vibSpring] (-2.5,-1.0) -- (-0.5,-1.0);
    \node at (-1.5,1.5) {$k_1$};

    \draw[blue,thick] (-0.5,-1.5) rectangle (0.5,1.5);
    \node at (0,0) {$m_1$};

    \draw[thick,-Stealth] (0.8,0) -- (1.8,0) node[right] {$F_0\sin\omega t$};

    \draw[red,vibSpring] (0.5,1.2) -- (2.5,1.8);
    \draw[red,thick] (2.5,1.2) -- (2.5,2.2);
    \node[red] at (2.0,2.0) {$k_2$};

    \draw[red,thick] (2.0,2.5) rectangle (3.0,3.5);
    \node[red] at (2.5,3.0) {$m_2$};

    \node[blue] at (0,0) {};

    \draw[red,thick,dashed] (2.2,1.3) -- (3.3,1.3);
    \node[red] at (3.8,1.3) {$c_2$};
\end{tikzpicture}
\caption{动力吸振器（DVA/TMD）模型。主质量 $m_1$ 弹性连接于基础并受简谐力激励，副质量 $m_2$ 通过弹簧 $k_2$ 和阻尼 $c_2$ 连接，构成有阻尼的动力吸振器。}
\label{fig:ch12_tmd}
\end{figure}

\subsection{多自由度吸振器与分布式吸振}
\begin{itemize}
    \item \textbf{多频吸振器}：在主系统上并联多个调谐至不同频率的吸振器，可覆盖多个共振峰
    \item \textbf{连续型吸振器}：利用梁、板的模态作为吸振器，如约束阻尼梁
    \item \textbf{非线性吸振器}（Nonlinear Energy Sink, NES）：利用立方非线性弹簧，可实现宽频的能量靶向传递
\end{itemize}

\section{阻尼技术}
\subsection{自由阻尼层与约束阻尼层}
\textbf{自由阻尼层}（Free-layer damping）：将黏弹性材料直接粘贴在结构表面，其耗能来自材料的拉伸/压缩变形。厚度越大、材料损耗因子（$\eta$）越大，阻尼效果越好。适用于薄壁结构。

\textbf{约束阻尼层}（Constrained-layer damping）：在黏弹性层外覆加刚性约束层（如铝板），迫使黏弹性层发生剪切变形。约束层阻尼的耗能效率远高于自由阻尼层，尤其适用于低频域的振动控制。

\begin{myformula}
约束阻尼梁的有效损耗因子（Mead-Markus 理论）：
\[
\eta_{\text{eff}} = \eta\frac{XY}{1 + (2 + Y)X + (1 + Y)(1 + \eta^{2})X^{2}}
\]
其中 $X$ 为剪切参数（含黏弹性层剪切模量、厚度、频率），$Y$ 为几何参数（含约束层与基层刚度比）。
\end{myformula}

\subsection{黏弹性阻尼材料}
材料的损耗因子 $\eta$（令复模量 $E^{*} = E(1 + i\eta)$）为振动能量的衡量指标。典型黏弹性材料的 $\eta$ 在 $0.1 \sim 2.0$ 之间（金属仅为 $10^{-4} \sim 10^{-3}$）。

温度-频率等效原理（WLF 方程）：黏弹性材料可在频率和温度之间实现等效转换，设计者可通过调节成分来匹配所需工作温度和频率范围。

\subsection{复合结构阻尼}
将阻尼材料嵌入结构内部，形成三明治结构。复合材料层合板可通过对铺层角度和材料的合理设计，在特定模态引入较高阻尼。

\section{振动主动控制简介}
\subsection{前馈控制（Feedforward）}
检测与激励相关的参考信号，通过自适应滤波器生成反相控制信号。对于周期性激励（如旋转机械、螺旋桨噪声），前馈控制能实现窄带噪声/振动的极高效抑制。

代表性算法：Filtered-X LMS（Least Mean Square），其更新律为：
\[
\mathbf{w}(n+1) = \mathbf{w}(n) + 2\mu e(n)\,\hat{\mathbf{r}}(n)
\]
其中 $\mathbf{w}$ 为 FIR 滤波器系数，$\hat{\mathbf{r}}(n)$ 为经次级通道模型滤波的参考信号，$e(n)$ 为误差传感器信号。

\subsection{反馈控制（Feedback）}
以振动传感器信号 $\mathbf{y}(t)$ 为反馈，计算控制力 $\mathbf{u}(t) = -\mathbf{K}\mathbf{x}$。通过合理的传感器/作动器布局实现全局抑振。

\begin{myformula}
设系统状态方程为 $\dot{\mathbf{x}} = \mathbf{A}\mathbf{x} + \mathbf{B}\mathbf{u}$，线性二次型调节器（LQR）的最优控制为：
\[
\mathbf{u}(t) = -\mathbf{R}^{-1}\mathbf{B}^{T}\mathbf{P}\,\mathbf{x}(t)
\]
其中 $\mathbf{P}$ 满足代数 Riccati 方程：
\[
\mathbf{A}^{T}\mathbf{P} + \mathbf{P}\mathbf{A} - \mathbf{P}\mathbf{B}\mathbf{R}^{-1}\mathbf{B}^{T}\mathbf{P} + \mathbf{Q} = \mathbf{0}
\]
$\mathbf{Q}$ 和 $\mathbf{R}$ 分别为状态和控制能量代价的加权矩阵。
\end{myformula}

\subsection{压电作动器与传感器}
压电材料的正/逆压电效应可直接将机械应变与电信号耦合：
\begin{itemize}
    \item \textbf{PZT（锆钛酸铅）}：高机电耦合系数，适合作为高效作动器
    \item \textbf{PVDF（聚偏氟乙烯）}：柔性好，适合粘贴在曲面作传感器
    \item \textbf{MFC（宏纤维复合材料）}：由纤维压电材料构成，兼具柔顺性和较大驱动力
\end{itemize}

\subsection{PID 振动控制}
PID 是最简单通用的反馈控制律：
\[
u(t) = K_p e(t) + K_i\int_0^{t} e(\tau)d\tau + K_d\frac{de(t)}{dt}
\]
其中 $e(t) = y_{\text{ref}} - y(t)$ 为误差信号。在振动控制中通常令 $y_{\text{ref}} = 0$（零位自动归零）。PID 增益通过试凑或 Ziegler-Nichols 法整定：

\begin{remark}
对于柔性结构的振动主动控制，PID 过高增益可能触发溢出（spillover）失稳——即控制力激发控制带宽外的高阶模态导致闭环不稳定。这是柔性结构主动控制必须关注的稳定性难题。
\end{remark}

\section{工程案例分析}
\subsection{建筑物抗震——TMD 调谐质量阻尼器}
台北 101 大楼（508\,\text{m}）在第 87--92 层悬挂直径 $5.5\,\text{m}$ 的 660 吨钢球，形成全球最著名的 TMD 装置。

\begin{myformula}
TMD 的最优调谐（$H_\infty$ 准则）：
\[
f_{\text{TMD}} = f_{\text{结构}}\frac{1}{1+\mu_{\text{eff}}}
\]
\[
\zeta_{\text{TMD,opt}} = \sqrt{\frac{3\mu_{\text{eff}}}{8(1+\mu_{\text{eff}})}}
\]
其中有效质量比 $\mu_{\text{eff}}$ 为 TMD 质量与结构广义质量之比（一般 $\mu \approx 0.01 \sim 0.05$）。
\end{myformula}

101 大楼 TMD 的参数实例：
\begin{itemize}
    \item 摆长 $\approx 11.5\,\text{m}$，调谐至结构一阶模态频率 $\approx 0.15\,\text{Hz}$（周期约 $6.7\,\text{s}$）
    \item 设置 8 个油压阻尼器提供附加阻尼
    \item 台风季中振幅可达 $\pm 1\,\text{m}$;能削减楼层加速度约 30\%--40\%
\end{itemize}

\begin{figure}[H]
\centering
\begin{tikzpicture}[scale=1.0]
    \draw[gray!30] (0,0) -- (2,0);
    \draw[pattern=north east lines] (0.5,-0.2) rectangle (1.5,0);

    \draw[thick] (1,-0.5) -- (1,-5.5);
    \draw[thick] (1,-5.5) -- (-1.5,-5.5);

    \draw[blue,thick] (-1.8,-5) rectangle (-0.2,-3.5);
    \node[blue] at (-1.0,-4.25) {$m_{\text{building}}$};

    \draw[vibSpring] (-3.5,-5) -- (-2.0,-5);
    \node at (-2.8,-4.7) {$k_1$};

    \draw[red,thick] (0.8,-1.5) -- (-0.5,-4.5);
    \draw[red,very thick,fill=red!20] (-0.8,-2.5) circle (0.6);
    \node[red] at (-0.8,-2.5) {\small $m_d$};

    \draw[red,vibSpring] (-2.0,-1.0) -- (0.5,-1.0);
    \draw[red,thick] (-2.0,-1.0) -- (-2.5,-1.5);

    \draw[gray,thick] (-3.5,-5) -- (-3.5,-4);
    \draw[gray,thick] (-3.5,-4) -- (-1.8,-4);

    \node[red] at (-1.5,-1.2) {$k_d$};
    \node at (-2.8,-3.5) {TMD};

    \draw[-Stealth,thick] (1.8,-5) -- (2.8,-5) node[right] {$F(t)$};
    \draw[gray] (2.8,-3.5) node {地震/风激励};
\end{tikzpicture}
\caption{台北101大楼 TMD（调谐质量阻尼器）的简化力学模型。$m_d$ 为调谐质量（660t），$k_d$ 和 $c_d$ 分别为等效摆刚度和阻尼，$m_1$ 和 $k_1$ 为建筑物广义质量和等效刚度。}
\label{fig:ch12_101_tmd}
\end{figure}

\subsection{旋转机械的转子动力学基础}
\subsubsection{Jeffcott 转子模型}
Jeffcott 转子（也称为 Laval 转子）是转子动力学的基础模型：质量为 $m$ 的刚性圆盘装在无质量弹性轴上，盘的质心偏离几何中心 $\varepsilon$（偏心距）。

\begin{myformula}
在固定坐标系 $(x, y)$ 中，假设各向同性轴（等刚度 $k = k_x = k_y$）：
\begin{align}
m\ddot{x} + c\dot{x} + kx &= m\varepsilon\Omega^{2}\cos\Omega t + mg \\
m\ddot{y} + c\dot{y} + ky &= m\varepsilon\Omega^{2}\sin\Omega t
\end{align}
其中 $\Omega$ 为轴转速。$y$ 方向稳态位移幅值为：
$$|Y| = \frac{m\varepsilon\Omega^{2}}{\sqrt{(k - m\Omega^{2})^{2} + (c\Omega)^{2}}}$$
\end{myformula}

\subsubsection{临界转速与不平衡响应}
当 $\Omega$ 接近横向固有频率 $\omega_n = \sqrt{k/m}$ 时，振幅急剧增大——此转速即为\textbf{临界转速}。在实际工程中：
\begin{itemize}
    \item 转子启动应迅速通过临界转速，避免在临界转速附近长时间驻留
    \item 平衡校正（动平衡）的目的是将残余偏心距 $\varepsilon$ 压至容许值以下
    \item 超过临界转速后转子呈自定心行为：质心位移趋于零而几何中心绕质心旋转
\end{itemize}

\subsection{汽车悬挂系统的振动分析——1/4 车模型}
1/4 车模型（Quarter-car Model）是汽车垂直振动分析的基本模型：
\begin{itemize}
    \item \textbf{悬挂质量 $m_s$}：1/4 车身质量
    \item \textbf{非悬挂质量 $m_u$}：轮组/车轴的等效质量
    \item \textbf{悬挂弹簧与阻尼器 $k_s, c_s$}：连接车身与车轴
    \item \textbf{轮胎刚度 $k_t$ 与轮胎阻尼 $c_t$}：连接车轴与路面
\end{itemize}

\begin{figure}[H]
\centering
\begin{tikzpicture}[scale=1.1]
    \draw[gray!30] (-2,-4.5) -- (2,-4.5);

    \draw[blue,thick] (-1.0,-3.5) rectangle (1.0,-2.5);
    \node[blue] at (0,-3.0) {$m_s$（悬挂质量）};

    \draw[vibSpring] (-1.0,-2.5) -- (-1.0,-1.5);
    \draw[thick] (-1.0,-2.5) -- (-1.8,-2.5);
    \draw[thick] (-1.8,-2.5) -- (-1.8,-1.5);
    \node at (-1.5,-1.8) {$k_s$};

    \draw[thick] (-0.2,-2.4) -- (-0.2,-1.6);
    \node at (0.3,-2.0) {$c_s$};

    \draw[red,thick] (-0.7,-1.5) rectangle (0.7,-0.5);
    \node[red] at (0,-1.0) {$m_u$（非悬挂质量）};

    \draw[vibSpring] (0.7,-0.5) -- (0.7,0.5);
    \draw[red,thick] (0.7,0.5) -- (0.7,1.0);
    \node at (1.2,0.2) {$k_t$};

    \draw[thick] (0.0,-0.4) -- (0.0,0.6);
    \node at (0.5,0.2) {$c_t$};

    \draw[green!50!black,very thick] (-2,1.0) -- (2,1.0);
    \node[green!50!black] at (0,1.5) {路面不平度 $y_r(t)$};
\end{tikzpicture}
\caption{1/4 车模型。$m_s$ 为悬挂质量（车身），$m_u$ 为非悬挂质量（车轮+车轴），$k_s$ 和 $c_s$ 为悬挂弹簧和阻尼器，$k_t$ 为轮胎刚度。路面不平度 $y_r(t)$ 作为基础位移激励输入。}
\label{fig:ch12_quarter_car}
\end{figure}

\begin{myformula}
1/4 车模型的运动方程：
\begin{align}
m_s\ddot{x}_s + c_s(\dot{x}_s - \dot{x}_u) + k_s(x_s - x_u) &= 0 \\
m_u\ddot{x}_u + c_s(\dot{x}_u - \dot{x}_s) + k_s(x_u - x_s) + k_t(x_u - y_r) + c_t(\dot{x}_u - \dot{y}_r) &= 0
\end{align}
悬挂设计的两个评价准则：（1）车身加速度最小（乘坐舒适性）;（2）轮胎动态载荷最小（行驶安全性）。
\end{myformula}

\subsection{桥梁的风致振动——Tacoma Narrows 的教训}
1940年11月7日，华盛顿州Tacoma Narrows大桥在风速约 $19\,\text{m/s}$ 时发生剧烈的扭转颤振，最终结构坍塌。这次事故彻底改变了桥梁风工程学的认知。

\textbf{气动弹性力学机理}：风速超过临界风速后，桥梁主梁的气动自激力与结构模态耦合产生负阻尼效应——系统的有效阻尼由正变负，能量从气流中持续注入结构，振幅指数增长，最终超过强度极限。

\textbf{现代抗风设计对策}：
\begin{itemize}
    \item 计算流体动力学（CFD）与风洞试验相结合，确定临界风速
    \item 主梁断面气动优化（加风嘴 $\to$ 流线型），涡流发生器
    \item 机翼型截面的高颤振临界风速
    \item TMD／调频液柱阻尼器（TLCD）附加阻尼
    \item 涡激振动（VIV）的振幅预测与规避设计（Scruton 数）
\end{itemize}

\section{Python 代码实现}
\subsection{隔振传递率计算与绘图}
\begin{mycode}{力传递率曲线绘制}
\begin{lstlisting}[language=Python]
import numpy as np
import matplotlib.pyplot as plt

def transmissibility(r, zeta):
    return np.sqrt(1 + (2 * zeta * r)**2) / \
           np.sqrt((1 - r**2)**2 + (2 * zeta * r)**2)

r = np.linspace(0, 5, 500)
zeta_vals = [0.0, 0.05, 0.1, 0.2, 0.5, 1.0]

plt.figure(figsize=(8, 5))
for z in zeta_vals:
    T = transmissibility(r, z)
    plt.plot(r, T, lw=1.5, label=f'$\\zeta={z}$')
plt.axhline(1, color='gray', ls='--', lw=0.8)
plt.axvline(np.sqrt(2), color='gray', ls='--', lw=0.8)
plt.text(np.sqrt(2)+0.1, 0.15, '$\\sqrt{2}$', fontsize=12)
plt.xlabel('$\\omega/\\omega_n$'); plt.ylabel('$T_f$')
plt.title('力传递率曲线'); plt.legend(); plt.grid(alpha=0.3)
plt.ylim(0, 5); plt.tight_layout(); plt.show()
\end{lstlisting}
\end{mycode}

\subsection{TMD 参数优化}
\begin{mycode}{Den Hartog 最优 TMD 参数计算与幅频响应}
\begin{lstlisting}[language=Python]
import numpy as np
import matplotlib.pyplot as plt

def tmd_optimal(mu):
    alpha_opt = 1.0 / (1 + mu)
    zeta_opt = np.sqrt(3 * mu / (8 * (1 + mu)))
    return alpha_opt, zeta_opt

def tmd_amplification(r, mu, alpha, zeta_a, zeta1=0.0):
    A = 1 - r**2
    B = 2 * zeta_a * r
    C = mu * r**2
    D = (alpha**2 - r**2)
    E = 2 * zeta_a * r
    F = (1 - r**2 - mu * r**2)
    G = 2 * zeta1 * r

    num = np.sqrt(A**2 + B**2)
    den = np.sqrt((A*D - B*E - C*(1 - alpha**2))**2 + (A*E + B*D - 2*C*zeta_a*r)**2)
    return num / den

mu_vals = [0.01, 0.05, 0.1]
r = np.linspace(0.5, 1.5, 500)
plt.figure(figsize=(8, 5))
for mu in mu_vals:
    alpha, zeta = tmd_optimal(mu)
    X1 = tmd_amplification(r, mu, alpha, zeta)
    plt.plot(r, X1, lw=1.5, label=f'$\\mu={mu}, \\alpha={alpha:.3f}, \\zeta={zeta:.3f}$')
plt.axvline(1, color='gray', ls='--', lw=0.8)
plt.xlabel('$\\omega/\\omega_1$'); plt.ylabel('$|X_1| / (F_0/k_1)$')
plt.title('有阻尼 TMD 幅频响应（Den Hartog 最优）')
plt.legend(); plt.grid(alpha=0.3); plt.tight_layout(); plt.show()
\end{lstlisting}
\end{mycode}

\section{习题与思考}
\begin{enumerate}
    \item 推导单自由度力隔振系统的力传递率公式，证明所有阻尼比对应的传递率曲线均在 $\omega/\omega_n = \sqrt{2}$ 处交汇于 $T_f = 1$。
    \item 设计一种无阻尼动力吸振器，使主系统在 $\omega = 15\,\text{rad/s}$ 处位移完全消除。已知 $m_1 = 10\,\text{kg}$, $k_1 = 2000\,\text{N/m}$，确定 $m_2$ 和 $k_2$。评估吸振器引入后系统的稳态副质量位移。
    \item 分别使用 $H_\infty$ 和 $H_2$ 优化准则，对质量比 $\mu = 0.05$ 的 TMD 进行参数设计，画出两条幅频曲线并讨论准则差异及其物理合理性。
    \item 推导 Jeffcott 转子在临界转速前后的相位差（放大因子），解释何谓"自定心"效应。
    \item 分析 1/4 车辆模型中悬挂刚度和阻尼对乘坐舒适性（ISO 2631 标准）的影响，绘制舒适性-悬挂刚度曲线。
    \item 探讨无阻尼吸振器的局限性：当激励频率偏离吸振器调谐频率时，主系统振幅会如何增长？为什么需要引入阻尼？
    \item 讨论主动控制中的"溢出"（spillover）问题及其对柔性结构控制稳定性的影响，提出可能的缓解方案。
    \item 编程实现 Jeffcott 转子的不平衡响应计算，画出幅频曲线和轴心轨迹。
\end{enumerate}
